% Source: https://artofproblemsolving.com/community/c3h1761161_partial_fraction
% Stable identity: aops-topic-1761161-q01
Multiplying by $n(n+2)$ gives $1=A(n+2)+Bn=(A+B)n+2A$ for all $n$. Comparing coefficients, $2A=1$ and $A+B=0$, so $\boxed{A=\dfrac{1}{2},\ B=-\dfrac{1}{2}}$. Indeed $\dfrac{1}{2n}-\dfrac{1}{2(n+2)}=\dfrac{(n+2)-n}{2n(n+2)}=\dfrac{1}{n(n+2)}$.
